% Part: proof-theory
% Chapter: normalization
% Section: normalization-thm

\documentclass[../../../include/open-logic-section]{subfiles}

\begin{document}

\olfileid{pt}{nor}{thm}

\olsection{স্বাভাবিকীকরণ উপপাদ্য}

!!^a{proof}-এ কোনো কর্তনখণ্ড না থাকলে
তা \emph{স্বাভাবিক} রূপে আছে। স্পষ্টত,
এর সমতুল্য কথা হলো প্রমাণটিতে আর
বিন্যাস-রূপান্তর বা হ্রাস-রূপান্তর
কোনোটিই প্রয়োগ করা যায় না।
!!a{proof} \emph{স্বাভাবিকীকরণ}
মানে বিন্যাস ও হ্রাস রূপান্তর
প্রয়োগ করে সেটিকে কর্তনহীন স্বাভাবিক
প্রমাণে রূপ দেওয়া। এটি সব সময়
সম্ভব—এই হলো \emph{স্বাভাবিকীকরণ
উপপাদ্য}।

\begin{thm}\ollabel{thm:normalization}
প্রত্যেক \Log{N1i}-!!{proof}~$\delta$
স্বাভাবিকীকৃত হয়; অর্থাৎ কর্তনহীন
!!a{proof}~$\delta'$-তে রূপান্তরিত
করা যায়।
\end{thm}

\begin{proof}
  $\delta$-র কর্তনমাত্রা ও
  কর্তনদৈর্ঘ্যের উপর যুগ্ম আরোহ।
  আরোহের অনুমান: কম কর্তনমাত্রার,
  অথবা একই কর্তনমাত্রায় কম
  কর্তনদৈর্ঘ্যের যেকোনো !!{proof}
  স্বাভাবিকীকৃত হয়।

  কর্তনদৈর্ঘ্য~$0$ হলে কোনো কর্তন
  নেই; $\delta$ আগেই স্বাভাবিক।
  তাই ধরি, $\delta$-র কর্তনমাত্রা
  ও কর্তনদৈর্ঘ্য উভয়ই~$> 0$।
  প্রমাণবৃক্ষের সবচেয়ে উপরের,
  তার মধ্যে সবচেয়ে ডানের,
  সর্বাধিক কর্তনখণ্ডটি বেছে নিই।
  খণ্ডটির দৈর্ঘ্য~$>1$ হলে
  বিন্যাস-রূপান্তর প্রয়োগ করি।
  দৈর্ঘ্য~$1$ হলে সংশ্লিষ্ট
  হ্রাস-রূপান্তর প্রয়োগ করি।
  এতে এমন নতুন !!{proof} পাই
  যার কর্তনমাত্রা একই কিন্তু
  কর্তনদৈর্ঘ্য কম, অথবা
  কর্তনমাত্রাই কম; আরোহের
  অনুমান তখন প্রযোজ্য।
\end{proof}

উপরের প্রমাণে সব সময় সবচেয়ে
উপরের, তার মধ্যে সবচেয়ে ডানের
সর্বাধিক কর্তন
নেওয়ার কৌশলটি রূপান্তরগুলি কোন
ক্রমে প্রয়োগ করতে হবে তার একটি
নির্দিষ্ট পথ দেখায়, যার ফলে
!!a{proof} স্বাভাবিক !!{proof}-এ
পরিণত হয়। আশ্চর্যের বিষয়,
ক্রমটি আসলে গুরুত্বপূর্ণ নয়:
বিন্যাস ও হ্রাস রূপান্তর যে
ক্রমেই প্রয়োগ করি না কেন,
শেষ পর্যন্ত স্বাভাবিক প্রমাণ
পাওয়া যায়।

\Log{N2i}-এর জন্যও খণ্ড ও
কর্তনের সংজ্ঞা সহজে দেওয়া যায়।
আর \Log{N1i}-র বিন্যাস ও হ্রাস
রূপান্তর সরাসরি \Log{N2i}-তে
অনুবাদ করা যায়। ফলে
\Log{N2i}-র জন্যও
স্বাভাবিকীকরণ উপপাদ্য চলে।

\begin{thm}\ollabel{thm:normalization-N2i}
প্রত্যেক \Log{N2i}-!!{proof}~$\delta$
স্বাভাবিকীকৃত হয়; অর্থাৎ কর্তনহীন
!!a{proof}~$\delta'$-তে রূপান্তরিত
করা যায়।
\end{thm}

\end{document}
